\chapter{Various}

\section{Game Theory}
	Impartial games: Both players have the same set of moves from any position.
	Normal play: Last move wins
	\subsection{Grundy Numbers}
	For some position $P$:    
	\[
		G(P) = \mathrm{mex}\{\, G(P') \mid P' \text{ is reachable from } P \,\} 
	\]
	Idea: $G(P) = 0$ iff $P$ is losing. \\
	Sprague-Grundy theorem: if you combine independent games $P_1, P_2, \dots, P_k$: 
	\[
		G(\text{sum}) = G(P_1) \oplus G(P_2) \oplus \dots \oplus G(P_k)
	\]
	Green Hackenbush: a graph with some vertices on the ground; a move deletes an edge, then everything no longer connected to the ground disappears; last move wins. \\
	A vertex's subtrees can be replaced by a single line of length $\bigoplus_{c \in \text{children}} (1 + \text{line}(c))$ (the connecting edge counts). A cycle can be contracted into a single node with its other edges sticking out; each cycle edge becomes a loop, worth a stalk of length 1. All ground vertices count as one. 

	\subsection{Nim variants}
	\textbf{Misere nim}: normal nim, unless all piles are $\le1$: then the first player wins iff the number of 1-piles is even.
	\textbf{Staircase nim} (coins move from step $i$ to $i-1$, step 0 = ground): only odd steps count, nim on their XOR.
	\textbf{Moore's Nim$_k$} (take from up to $k$ piles): losing iff for every bit the number of piles with that bit set is divisible by $k+1$.
	\textbf{Wythoff} (take from one pile, or equally from both): P-positions $(\lfloor m\varphi\rfloor,\lfloor m\varphi\rfloor+m)$, $\varphi=\frac{1+\sqrt5}2$: $(0,0),(1,2),(3,5),(4,7),(6,10),\dots$
	\textbf{Lasker's nim} (take, or split a pile in two): $G(4m{+}1)=4m{+}1$, $G(4m{+}2)=4m{+}2$, $G(4m{+}3)=4m{+}4$, $G(4m{+}4)=4m{+}3$.

\section{Intervals}
	\kactlimport{IntervalContainer.h}
	\kactlimport{IntervalCover.h}
	\kactlimport{ConstantIntervals.h}

\section{Misc. algorithms}
	\kactlimport{TernarySearch.h}
	\kactlimport{CyclicMax.h}
	\kactlimport{LIS.h} % TODO rewrite
	\kactlimport{FastKnapsack.h}

\section{Dynamic programming}
	\subsection{DP optimisations}
	\begin{itemize}[noitemsep]
	\item $dp_i=\min_{j<i}dp_j+w(j,i)$ with $w$ Monge ($w(a,c)+w(b,d)\le w(a,d)+w(b,c)$ for $a\le b\le c\le d$, e.g.\ $g(S_i-S_j)$, $g$ convex, $S$ sorted): argmin is monotone. CDQ D\&C $O(n\log^2n)$, or a monotone deque with binary search $O(n\log n)$.
	\item $dp_i=\min_jdp_j+b_ja_i$: CHT / Li Chao.
	\item Layers $dp_{t,i}=\min_jdp_{t-1,j}+w(j,i)$: D\&C DP $O(kn\log n)$.
	\end{itemize}
	\textbf{Aliens trick}: $f(k)$ = optimum with exactly $k$ items, $f$ convex. Solve the relaxed DP $g(\lambda)=\min_k\,(f(k)+\lambda k)$, returning the item count $c(\lambda)$ with ties broken to the smallest. Binary search the smallest integer $\lambda$ with $c(\lambda)\le k$ (range $\pm\max|f(i+1)-f(i)|$); then $f(k)=g(\lambda)-\lambda k$.

	\kactlimport{KnuthDP.h}
	\kactlimport{DivideAndConquerDP.h}

\section{Debugging tricks}
	\begin{itemize}
		\item \verb@signal(SIGSEGV, [](int) { _Exit(0); });@ converts segfaults into Wrong Answers.
			Similarly one can catch SIGABRT (assertion failures) and SIGFPE (zero divisions).
			\verb@_GLIBCXX_DEBUG@ failures generate SIGABRT (or SIGSEGV on gcc 5.4.0 apparently).
		\item \verb@feenableexcept(29);@ kills the program on NaNs (\texttt 1), 0-divs (\texttt 4), infinities (\texttt 8) and underflows (\texttt{16}).
	\end{itemize}
	\kactlimport{LeakSanitizer.h}

\section{Optimization tricks}
	\verb@__builtin_ia32_ldmxcsr(40896);@ disables denormals (which make floats 20x slower near their minimum value).
	\subsection{Bit hacks}
		\begin{itemize}
			\item \verb@x & -x@ is the least bit in \texttt{x}.
			\item \verb@for (int x = m; x; ) { --x &= m; ... }@ loops over all subset masks of \texttt{m} (except \texttt{m} itself).
			\item \verb@c = x & -x, r = x + c; (((r ^ x) >> 2) / c) | r@ is the next number after \texttt{x} with the same number of bits set.
			\item \verb@F0R (b, K) F0R (i, (1 << K))@ \\ \verb@  if (i & 1 << b) D[i] += D[i ^ (1 << b)];@ computes all sums of subsets.
		\end{itemize}
	\subsection{Pragmas}
		\begin{itemize}
			\item \lstinline{#pragma GCC optimize ("Ofast")} will make GCC auto-vectorize loops and optimizes floating points better.
			\item \lstinline{#pragma GCC target ("avx2")} can double performance of vectorized code, but causes crashes on old machines.
			\item \lstinline{#pragma GCC optimize ("O0", "trapv")} kills the program on integer overflows (but is really slow).
		\end{itemize}
	\kactlimport{FastMod.h}
	\kactlimport{FastInput.h}
	\kactlimport{BumpAllocator.h}
	\kactlimport{SmallPtr.h}
	\kactlimport{BumpAllocatorSTL.h}
	\kactlimport{Unrolling.h}

\section{Bigger types}
	\begin{itemize}[noitemsep]
		\item \texttt{\_\_int128} / \texttt{unsigned \_\_int128}: 128-bit ints, up to $\pm 1.7\times10^{38}$. No \texttt{cin/cout}: cast to \texttt{ll} or print digit by digit. Multiply is a few instructions, divide is slow (\tilde{}100 cycles). Use for exact $64\times64$-bit products, mulmod, sums up to $10^{36}$.
		\item \texttt{long double}: 80-bit x87 on Linux (64-bit mantissa, \tilde{}19 digits), \tilde{}3$\times$ slower than \texttt{double}. Use \texttt{sqrtl}, \texttt{powl}, \texttt{fabsl}, \texttt{\%Lf}.
		\item \texttt{\_\_float128}: 113-bit mantissa (\tilde{}34 digits), software emulated (libgcc calls, \tilde{}25$\times$ slower than \texttt{double}). Arithmetic and casts just work; \texttt{sqrtq}, \texttt{expq} and printing need \texttt{<quadmath.h>} + \texttt{-lquadmath}, so print via \texttt{(long double)} or peel digits manually.
	\end{itemize}
	\kactlimport{BigInt.h}

\section{Other}
	\kactlimport{GrayCode.h}
	% \kactlimport{Timer.h}
	% \kactlimport{SIMD.h}
	\kactlimport{STL.h}
	% \kactlimport{Date.h}
	\kactlimport{Python.py}
\section{Test Session}
	\kactlimport{TestSession.h}